% =====================================================================
%  Redes Neuronales 2026 (FaMAF-UNC) - Trabajo Práctico 1
%  Boceto de informe: "El modelo de neuronas de Izhikevich"
%
%  Cómo usarlo:
%   - Overleaf: New Project > Upload Project > subí el .zip con main.tex y
%     la carpeta figuras/. Compilá con pdfLaTeX (el botón "Recompile").
%   - Local: latexmk -pdf main.tex   (o pdflatex main.tex dos veces).
%
%  Todo lo que está entre [corchetes] o marcado con TODO es para revisar o
%  reemplazar. Los números del texto salen de la guía paso a paso del TP
%  (carpeta tp1-izhikevich/): verifíquenlos contra sus propias corridas.
%  Límite de la cátedra: el PDF final NO debe superar las 4 páginas.
% =====================================================================
\documentclass[10pt,a4paper,twocolumn]{article}

% ---------- Idioma y codificación ----------
\usepackage[utf8]{inputenc}
\usepackage[T1]{fontenc}
\usepackage[spanish,es-tabla,es-noshorthands]{babel} % "Tabla" en vez de "Cuadro"
\usepackage{lmodern}

% ---------- Página ----------
\usepackage[a4paper,margin=1.8cm,columnsep=0.7cm]{geometry}
\usepackage{microtype}

% ---------- Matemática y unidades ----------
\usepackage{amsmath,amssymb}
\usepackage{siunitx}
\sisetup{output-decimal-marker={,}}          % 0,04 en vez de 0.04
\DeclareSIUnit{\milivolt}{mV}

% ---------- Figuras y tablas ----------
\usepackage{graphicx}
\graphicspath{{figuras/}}                     % las imágenes van en figuras/
\usepackage{subcaption}                       % paneles (a), (b) dentro de una figura
\usepackage{booktabs}                         % líneas lindas en tablas
\usepackage[font=small,labelfont=bf]{caption}

% ---------- Enlaces ----------
\usepackage{xcolor}
\usepackage[colorlinks=true,linkcolor=blue!50!black,citecolor=blue!50!black,urlcolor=blue!50!black]{hyperref}

% ---------------------------------------------------------------------
% \figura{archivo}{ancho}: incluye la imagen si existe en figuras/; si
% todavía no la generaron, dibuja un recuadro con el nombre. Así el
% documento compila siempre.
% ---------------------------------------------------------------------
\newcommand{\figura}[2]{%
  \IfFileExists{figuras/#1}%
    {\includegraphics[width=#2]{#1}}%
    {\fbox{\parbox[c][0.45#2][c]{\dimexpr#2-2\fboxsep-2\fboxrule}{\centering\small\ttfamily figuras/#1\\(falta generar)}}}%
}

% =====================================================================
\title{\textbf{Simulación numérica del modelo de neuronas de Izhikevich:\\
tipos de disparo, error numérico y dinámica de red}}
\author{%
  [Nombre Apellido 1], [Nombre Apellido 2] y [Nombre Apellido 3]\\[2pt]
  \small Facultad de Matemática, Astronomía, Física y Computación, Universidad Nacional de Córdoba
}
\date{Redes Neuronales 2026 -- Trabajo Práctico 1}

\begin{document}

% El resumen ocupa todo el ancho, debajo del título. Las llaves {...} protegen
% los corchetes [ ] del texto, que si no cerrarían \twocolumn[...].
\twocolumn[{
  \begin{@twocolumnfalse}
    \maketitle
    \begin{abstract}
    % 5 a 10 líneas: objetivo + metodología + resultados principales.
    Estudiamos numéricamente el modelo de neuronas de Izhikevich, un sistema
    bidimensional de ecuaciones diferenciales con reseteo que reproduce la variedad
    de patrones de disparo de las neuronas corticales con un costo computacional
    mínimo. Integramos el modelo con el método de Runge--Kutta de orden 4
    ($h=\num{0.1}$\,ms) y reproducimos los ocho tipos de disparo de la figura~2 de
    la Ref.~\cite{izhikevich2003}: regular spiking, intrinsically bursting,
    chattering, fast spiking, tálamo-cortical, resonador y low-threshold spiking.
    Mostramos que la condición inicial en el reposo es crítica para los casos de
    umbral bajo, y que el reseteo discreto reduce a orden~1 la convergencia de RK4
    en los tiempos de disparo. [Si hacen la parte 2:] Finalmente, simulamos una red de
    1000 neuronas acopladas con ruido mediante Euler--Maruyama y observamos ritmos
    $\alpha$ y $\gamma$, que se atenúan al reemplazar la activación escalón por una
    función lineal $z(v)$.
    \end{abstract}
    \vspace{1em}
  \end{@twocolumnfalse}
}]

% =====================================================================
\section{Introducción}
% Qué es un modelo de neurona, el compromiso realismo/eficiencia, citas, objetivo.
Las neuronas se comunican mediante potenciales de acción o \emph{spikes}:
excursiones bruscas y breves del potencial de membrana. Los modelos de neurona
de disparo enfrentan un compromiso entre realismo y eficiencia. El modelo de
Hodgkin y Huxley~\cite{hodgkin1952} describe el fenómeno con gran detalle
biofísico, pero con varias ecuaciones acopladas y decenas de parámetros, lo
que lo vuelve costoso para simular redes grandes. En el otro extremo, el
modelo \emph{integrate-and-fire} es muy eficiente pero no reproduce la riqueza
de patrones de disparo observada experimentalmente.

Izhikevich~\cite{izhikevich2003} propuso un modelo que resuelve ese
compromiso: dos ecuaciones con un único término no lineal, obtenidas por
reducción a la forma normal de las bifurcaciones que dan origen al
disparo~\cite{izhikevich2007}. Con sólo cuatro parámetros reproduce los
principales tipos de disparo de neuronas corticales y talámicas, y permite
simular decenas de miles de neuronas en tiempo real.

El objetivo de este informe es integrar numéricamente el modelo, reproducir
sus tipos de disparo y analizar la precisión del método numérico
[y la dinámica colectiva de una red de neuronas acopladas]. % TODO

% =====================================================================
\section{Teoría}
\subsection{El modelo de Izhikevich}
El estado de la neurona se describe con el potencial de membrana $v(t)$, en mV,
y una variable de recuperación $u(t)$, que representa en forma efectiva la
activación de corrientes de K$^+$ y la inactivación de corrientes de Na$^+$, y
aporta una realimentación negativa sobre $v$. Con $t$ en ms,
\begin{align}
  \dot v &= g_2\,v^2 + g_1\,v + g_0 - u + I(t), \label{eq:v}\\
  \dot u &= b\,(c\,v - u), \label{eq:u}
\end{align}
donde $I(t)$ es la corriente de entrada (sináptica o inyectada). Cuando
$v\geq v_+$ se considera que la neurona disparó y se aplica el reseteo
\begin{equation}
  v \leftarrow v_-, \qquad u \leftarrow u + \Delta u .
  \label{eq:reset}
\end{equation}
Los parámetros (entre paréntesis, la notación de la Ref.~\cite{izhikevich2003}) son:
$b$ ($a$), la escala temporal de $u$: valores chicos implican recuperación lenta;
$c$ ($b$), la sensibilidad de $u$ a las fluctuaciones subumbrales de $v$;
$v_-$ ($c$), el potencial de reseteo; y $\Delta u$ ($d$), el salto de $u$ en cada
spike. Los coeficientes $g_2=\num{0.04}$, $g_1=5$ y $g_0=140$ fijan las unidades,
y $v_+=\SI{30}{\milivolt}$ es el pico del spike.

Para $I$ constante, el reposo es la raíz menor de
$g_2v^2+(g_1-c)v+g_0+I=0$, con $u^*=c\,v^*$; para $c=\num{0.2}$ e $I=0$ resulta
$v^*=\SI{-70}{\milivolt}$. [Si hay lugar: el reposo se pierde en
$I_\mathrm{Hopf}\approx\num{3.8}$ e $I_\mathrm{SN}=4$.] % TODO

\subsection{Método numérico}
Con $\mathbf{x}=(v,u)$ y $\dot{\mathbf{x}}=\mathbf{f}(t,\mathbf{x})$, usamos
Runge--Kutta de orden~4:
\begin{equation}
  \mathbf{x}_{n+1} = \mathbf{x}_n + \tfrac{h}{6}\left(\mathbf{k}_1 + 2\mathbf{k}_2 + 2\mathbf{k}_3 + \mathbf{k}_4\right),
\end{equation}
con $\mathbf{k}_1=\mathbf{f}(t_n,\mathbf{x}_n)$,
$\mathbf{k}_{2,3}=\mathbf{f}(t_n+\tfrac h2,\mathbf{x}_n+\tfrac h2\mathbf{k}_{1,2})$ y
$\mathbf{k}_4=\mathbf{f}(t_n+h,\mathbf{x}_n+h\mathbf{k}_3)$.
Después de cada paso, una función condicionante verifica $v_{n+1}\geq v_+$ y,
si se cumple, aplica el reseteo~\eqref{eq:reset}. Integramos en
$t\in[0,200]$\,ms con $h=\num{0.1}$\,ms, partiendo del reposo de cada caso,
$v(0)=v^*(c)$ y $u(0)=c\,v(0)$ (para $c=\num{0.2}$, $v(0)=\SI{-70}{\milivolt}$).

% =====================================================================
\section{Resultados}
% Regla: la leyenda dice QUÉ se grafica; el texto dice QUÉ SE VE; el porqué va en Discusión.

\begin{figure}[t]
  \centering
  \figura{f1_base.png}{\columnwidth}
  \caption{Evolución temporal de $v(t)$ (arriba), $u(t)$ (medio) e $I(t)$
  (abajo) para $b=\num{0.02}$, $c=\num{0.2}$, $v_-=\SI{-65}{\milivolt}$,
  $\Delta u=2$ y un escalón $I=10$ desde $t=10$\,ms, integrados con RK4 y
  $h=\num{0.1}$\,ms. En azul, los spikes igualados a $v_+$.}
  \label{fig:base}
\end{figure}

En la Figura~\ref{fig:base} se observa que, mientras $I=0$, la neurona permanece
en reposo en $\SI{-70}{\milivolt}$. Al encenderse el escalón emite el primer
spike a los \num{3.5}\,ms; los intervalos entre spikes (ISI) crecen de
\num{3.6} a \num{15.2}\,ms y desde $t\approx45$\,ms el disparo es periódico
con ISI $\approx19$\,ms ($\approx53$\,Hz). La variable $u$ salta $\Delta u$
en cada spike y relaja hacia $c\,v$ entre spikes. % TODO: verificar números

Para reproducir la figura~2 de la Ref.~\cite{izhikevich2003} usamos los
parámetros de la Tabla~\ref{tab:parametros} y las corrientes
$I_1$: 10 desde $t=25$; $I_2$: \num{2.5} desde $t=150$;
$I_3$: $-15$ hasta $t=100$ y 0 después; e
$I_4$: \num{0.2} desde $t=10$ con un pulso de 10 en $[100,101)$\,ms.

\begin{table}[t]
  \centering
  \caption{Parámetros de cada tipo de neurona. En todos los casos
  $g_2=\num{0.04}$, $g_1=5$, $g_0=140$ y $v_+=\SI{30}{\milivolt}$.}
  \label{tab:parametros}
  \small
  \begin{tabular}{lccccc}
    \toprule
    Caso & $b$ & $c$ & $v_-$ [mV] & $\Delta u$ & $I$ \\
    \midrule
    RS  & \num{0.02} & \num{0.2}  & $-65$ & 8          & $I_1$ \\
    IB  & \num{0.02} & \num{0.2}  & $-55$ & 4          & $I_1$ \\
    CH  & \num{0.02} & \num{0.2}  & $-50$ & 2          & $I_1$ \\
    FS  & \num{0.1}  & \num{0.2}  & $-65$ & 2          & $I_1$ \\
    TC1 & \num{0.02} & \num{0.25} & $-65$ & \num{0.05} & $I_2$ \\
    TC2 & \num{0.02} & \num{0.25} & $-65$ & \num{0.05} & $I_3$ \\
    RZ  & \num{0.1}  & \num{0.26} & $-65$ & 2          & $I_4$ \\
    LTS & \num{0.02} & \num{0.25} & $-65$ & 2          & $I_1$ \\
    \bottomrule
  \end{tabular}
\end{table}

% figure* = figura a doble columna (va arriba de la página siguiente).
\begin{figure*}[t]
  \centering
  \figura{f2_figura2.png}{0.72\textwidth}
  \caption{Potencial de membrana $v(t)$ (azul) y corriente de entrada $I(t)$
  (negro, escala arbitraria) para las parametrizaciones RS, IB, CH, FS, TC (dos
  protocolos), RZ y LTS de la Tabla~\ref{tab:parametros}, integradas con RK4 y
  $h=\num{0.1}$\,ms desde el reposo de cada caso. Spikes igualados a $v_+$.}
  \label{fig:tipos}
\end{figure*}

En la Figura~\ref{fig:tipos} se observa que RS produce un tren de spikes cuyo
ISI crece de 17 a 45\,ms, mientras que FS dispara a $\approx133$\,Hz sin
adaptación apreciable. IB emite una ráfaga inicial de tres spikes seguida de
spikes simples cada 31\,ms, y CH dispara ráfagas de cinco a ocho spikes
separadas por $\approx60$\,ms. En TC1 un escalón de \num{2.5} alcanza para un
disparo tónico, mientras que en TC2 la neurona se hiperpolariza hasta
$\approx\SI{-82}{\milivolt}$ y al liberar la corriente dispara una ráfaga de
rebote. En RZ el pulso de 1\,ms lleva a la neurona de una oscilación subumbral
amortiguada a un disparo sostenido, y LTS dispara a alta frecuencia con
adaptación notoria (ISI de 3 a \num{13.5}\,ms). % TODO: verificar números

\begin{figure}[t]
  \centering
  \figura{f6_error.png}{\columnwidth}
  \caption{Izquierda: máximo error absoluto en los tiempos de disparo de RS en
  función de $h$, para Euler, RK4 y RK4 con localización del cruce por
  interpolación lineal. Derecha: $v(t)$ con Euler ($h=\num{0.5}$\,ms) y RK4
  ($h=\num{0.1}$\,ms); las líneas punteadas marcan los disparos de referencia.}
  \label{fig:error}
\end{figure}

La Figura~\ref{fig:error} muestra que el error en los tiempos de disparo de
RK4 decrece linealmente con $h$, igual que el de Euler pero con una constante
$\approx\num{2.5}$ veces menor: para $h=\num{0.1}$\,ms es de \num{0.68}\,ms.
Al localizar el cruce de $v_+$ por interpolación el error baja a
\num{0.09}\,ms. % TODO

% ---- Segunda parte (opcional). Si no la hacen, borren esta subsección. ----
\subsection{Red de neuronas acopladas}
Simulamos una red de $n_e=800$ neuronas excitatorias y $n_i=200$
inhibitorias con el método de Euler--Maruyama, con la matriz de acoplamiento
$a_{ij}$ de la consigna. La Figura~\ref{fig:red} compara la activación
escalón $\Theta(v)$ con la función lineal $z(v)=(87+v)/450-\num{0.0193}$. Con
$\Theta$ aparecen bandas verticales de disparo sincronizado, como en la
figura~3 de la Ref.~\cite{izhikevich2003}; con $z(v)$ la tasa media es
similar ($\num{9.7}$ contra $\num{9.1}$\,Hz) pero las bandas casi desaparecen. % TODO

\begin{figure*}[t]
  \centering
  \begin{subfigure}{0.49\textwidth}\centering
    \figura{f8_red_theta.png}{\linewidth}\caption{Activación $\Theta(v)$}\label{fig:red-theta}
  \end{subfigure}\hfill
  \begin{subfigure}{0.49\textwidth}\centering
    \figura{f8_red_z.png}{\linewidth}\caption{Activación $z(v)$}\label{fig:red-z}
  \end{subfigure}
  \caption{Arriba: raster de disparos de la red de 1000 neuronas
  (Euler--Maruyama, $\tau=\num{0.5}$\,ms); las neuronas 1 a 800 son
  excitatorias y el resto inhibitorias. Abajo: $v(t)$ de la neurona~1. Ambas
  corridas usan la misma red y la misma semilla de ruido.}
  \label{fig:red}
\end{figure*}

% =====================================================================
\section{Discusión}
% ¿Se reproducen los tipos? Mecanismos. Fuentes de error. Importancia del modelo.
Las simulaciones reproducen cualitativamente los ocho tipos de disparo de la
Ref.~\cite{izhikevich2003}. Las diferencias cuantitativas (cantidad de spikes
por ráfaga, el mínimo de TC2) se deben a que el paper no especifica las
corrientes de entrada exactas. Los mecanismos se leen en los parámetros: un
$\Delta u$ grande produce adaptación (RS); un $v_-$ alto favorece las ráfagas
(IB, CH); un $b$ grande da recuperación rápida (FS); y un $c$ grande baja el
umbral (LTS, TC) y da oscilaciones subumbrales (RZ). Los casos de umbral bajo
son sensibles a la condición inicial: partiendo de $\SI{-70}{\milivolt}$ en lugar
del reposo aparecen spikes espurios antes del estímulo.

Como fuentes de error numérico destacamos: (i) el reseteo se detecta recién en
el primer punto de la grilla con $v\geq v_+$, lo que introduce un error
$\mathcal{O}(h)$ en cada tiempo de disparo, que se acumula de un spike al
siguiente y anula la ventaja del orden~4 de RK4 (Fig.~\ref{fig:error});
(ii) la rigidez del término $v^2$ cerca del spike, que con $h$ grande puede
hacer divergir $v$ dentro de un paso; (iii) el pico muestreado no alcanza
exactamente $v_+$; y (iv) en sistemas biestables como RZ, un error pequeño
puede cambiar de régimen la trayectoria.

[Si hacen la parte 2: interpretar por qué $z(v)$ atenúa la sincronización.]
Finalmente, el modelo de Izhikevich combina la riqueza dinámica de los modelos
biofísicos con la eficiencia de los modelos de integración y disparo, lo que lo
hace adecuado para simular redes grandes; la parte~2 muestra que de una red
sin plasticidad emergen ritmos corticales. % TODO

% =====================================================================
\section*{Agradecimientos}
% Opcional.
Agradecemos a [los docentes de la materia] por [...].

% =====================================================================
% Referencias "a mano". Alternativa con BibTeX: ver la guía.
\begin{thebibliography}{9}\small
\bibitem{izhikevich2003}
  E.~M. Izhikevich, ``Simple model of spiking neurons,''
  \emph{IEEE Trans.\ Neural Netw.} \textbf{14}(6), 1569--1572 (2003).

\bibitem{izhikevich2007}
  E.~M. Izhikevich, \emph{Dynamical Systems in Neuroscience: The Geometry of
  Excitability and Bursting}, MIT Press (2006).

\bibitem{hodgkin1952}
  A.~L. Hodgkin y A.~F. Huxley, ``A quantitative description of membrane
  current and its application to conduction and excitation in nerve,''
  \emph{J.\ Physiol.} \textbf{117}(4), 500--544 (1952).
\end{thebibliography}

\end{document}
